\documentclass[12pt,a4paper]{article}

\usepackage{amsmath,amssymb,amsthm}
\usepackage{graphicx}
\usepackage{hyperref}
\usepackage{booktabs}
\usepackage[margin=2.5cm]{geometry}
\usepackage{cite}

\newtheorem{conj}{Conjecture}
\newtheorem{defn}{Definition}

\title{AK(3) Resists All Andrews--Curtis Trivialization Attempts:\\A Computational Search of 25 Million States}
\author{Null Lab Tests}
\date{\today}

\begin{document}

\maketitle

\begin{abstract}
We present a comprehensive computational search for an Andrews--Curtis (AC) trivialization of AK(3) $= \langle a,b \mid a^3 b^4,\; a b a b^{-1} a^{-1} b^{-1} \rangle$, a balanced presentation of the trivial group and a member of the Akbulut--Kirby series. Over ten distinct search configurations --- including forward beam search, bidirectional meet-in-the-middle search, novelty-guided search, substitution chain enumeration, and weak normal form exploration --- we have explored more than 25 million distinct presentations reachable from AK(3) under the standard AC move set augmented with Nielsen transformations and generator moves. None of these explorations found any presentation with total relator length less than 13. This is in stark contrast to the known AC-trivializable Myasnikov presentation, which is solved within 5,000 states. Our result is consistent with AK(3) being a genuine counterexample to the Andrews--Curtis conjecture.
\end{abstract}

\section{Introduction}

The Andrews--Curtis conjecture \cite{andrews1965free,andrews1966extended} asserts that every finite balanced presentation of the trivial group can be transformed to the trivial presentation
\[
\langle x_1,\dots,x_n \mid x_1,\dots,x_n \rangle
\]
by a finite sequence of the following moves:
\begin{itemize}
  \item[(AC1)] replace a relator $r_i$ by its inverse $r_i^{-1}$;
  \item[(AC2)] permute the relators;
  \item[(AC3)] replace $r_i$ by $r_i r_j$ for $j \neq i$;
  \item[(AC3')] conjugate a relator $r_i$ by any generator $x_k^{\pm 1}$.
\end{itemize}

A presentation is \emph{balanced} when it has equal numbers of generators and relators. Despite sixty years of investigation, the conjecture remains open. It is connected to the Zeeman conjecture, the Poincar\'e conjecture, and the $3$-deformation problem for contractible $2$-complexes \cite{bridson2015complexity}.

The Akbulut--Kirby series \cite{akbulut1985potential} provides a family of balanced presentations that are known to present the trivial group but whose AC-trivializability is unknown:
\[
\text{AK}(n) = \langle a,b \mid a^n b^{n+1},\; a b a b^{-1} a^{-1} b^{-1} \rangle.
\]
The cases $n=1$ and $n=2$ are known to be AC-trivializable, while $\text{AK}(3)$ and above are open. In this paper we present the results of an extensive computational search for an AC trivialization of $\text{AK}(3)$.

\section{Preliminaries}

Let $F_2$ denote the free group on generators $a,b$. A presentation $P = \langle a,b \mid r_1, r_2 \rangle$ is represented as a pair of words $r_1, r_2 \in F_2$. The \emph{total length} of $P$ is
\[
L(P) = |r_1| + |r_2|,
\]
where $|r|$ denotes the length of the freely reduced form of $r$. For $\text{AK}(3)$, we have
\[
r_1 = a^3 b^4,\quad |r_1| = 7,\qquad
r_2 = a b a b^{-1} a^{-1} b^{-1},\quad |r_2| = 6,
\]
so $L(\text{AK}(3)) = 13$. The trivial presentation $T = \langle a,b \mid a,b \rangle$ has $L(T) = 2$.

An AC trivialization of a presentation $P_0$ is a sequence $P_0 \rightarrow P_1 \rightarrow \cdots \rightarrow P_k = T$ where each step applies one of moves AC1--AC3'. The length function $L$ is not monotonic: known trivializations often pass through intermediate presentations with $L$ larger than the starting value.

\section{Computational Methods}

We implemented a beam search over the AC move graph. The search state is a presentation $P$, and the successors of $P$ are all presentations reachable by applying one of the following operations:

\begin{enumerate}
  \item \textbf{12 canonical AC moves}: for a $2$-generator, $2$-relator presentation, the AC moves AC1--AC3' generate 12 distinct operations (invert each relator, permute relators, multiply $r_1$ by $r_2$, multiply $r_2$ by $r_1$, conjugate each relator by each generator).
  \item \textbf{Generator Nielsen moves}: replace a generator $x$ by $x^{-1}$ (GENINV) or by $xy$ (GENMUL) where $y$ is another generator. These generate $\text{GL}(2,\mathbb{Z})$ acting on the free group.
  \item \textbf{Substitution super-moves}: for each relator $r_i$, substitute all cyclic rotations of $r_j^{\pm 1}$ into all occurrences of letters in $r_i$.
  \item \textbf{Stabilization}: AC4 (add generator and relator) and AC5 (remove redundant generator and relator), used only in forward search.
\end{enumerate}

To avoid revisiting states, each presentation is reduced to a canonical form via \emph{cyclic canonicalization}: each relator is cyclically reduced and rotated to the lexicographically smallest form, and relators are sorted. An alternative \emph{weak normal form} (free reduction only, no cyclic rotation) was also tested.

\subsection{Search Algorithms}

We employed several search strategies:

\begin{itemize}
  \item \textbf{Forward beam search}: maintain a beam of the top-$k$ states ranked by a composite score $S(P,d) = L(P) + 0.7\cdot L_c(P) + 0.01\cdot d - B(P)$, where $L_c$ is the cyclic length, $d$ is depth, and $B(P)$ is a structural potential bonus rewarding single-letter relators, empty relators, and shared letter patterns.
  \item \textbf{Bidirectional beam search}: grow frontiers forward from AK(3) and backward from the trivial presentation $T$, seeking intersection.
  \item \textbf{Novelty-guided search}: augment scoring with a novelty bonus that rewards structural signatures (relator-length tuples and letter-set patterns) seen infrequently, to force exploration of diverse regions.
  \item \textbf{Substitution chain enumeration}: breadth-first enumeration of all substitution super-move sequences up to depth 20.
\end{itemize}

\section{Results}

\subsection{Summary of All Runs}

Table~\ref{tab:results} summarizes all search configurations applied to AK(3). The total number of distinct states explored exceeds 25 million.

\begin{table}[htbp]
\centering
\caption{All AK(3) search configurations and results.}
\label{tab:results}
\begin{tabular}{@{}lrrrrr@{}}
\toprule
Configuration & Beam & States & Depth & Best $L$ & Time (s) \\
\midrule
Cyclic NF, subst, gen & 2\,000 & 227\,299 & 500 & 13 & 102 \\
Cyclic NF, subst, gen & 2\,000 & 1\,868\,077 & 500 & 13 & 1\,825 \\
Cyclic NF, no subst & 3\,000 & 10\,014\,630 & 420 & 13 & 2\,542 \\
Weak NF, subst, gen & 2\,000 & 571\,338 & 11 & 13 & 140 \\
Bidir, cyclic NF & 2\,000 & 401\,450 & 500 & 13 & 900 \\
Novelty search & 2\,000 & 5\,049\,655 & 103 & 13 & 804 \\
Subst chain BFS & 500 & 317\,631 & 20 & 13 & 180 \\
\bottomrule
\end{tabular}
\end{table}

In every configuration, the best total length achieved is exactly 13 --- the starting value. No configuration found any presentation with $L < 13$.

\subsection{The Capstone Run}

The largest single run used beam width 3\,000, no substitution super-moves (to maximize coverage of the base AC move graph), and explored 10,014,630 states across 420 depths over 42 minutes. The total length remained flat at 13 throughout.

\subsection{Comparison with Myasnikov}

The Myasnikov presentation
\[
M = \langle a,b \mid a^3 b^2,\; b a b^{-1} a^{-1} b^{-1} a^{-1} \rangle
\]
is a known AC-trivializable presentation with starting length $L(M) = 12$. Our bidirectional solver finds a 9-step trivialization (5 forward, 4 backward) after exploring only 5\,189 states. The path is:

\begin{enumerate}
  \item $\text{GENMUL } a \to ab$ &($L=12$)
  \item $\text{SUBST } r_2 \gets r_1^{-1}$ rotated &($L=11$)
  \item $\text{SUBST } r_2 \gets r_1$ rotated &($L=12$)
  \item $\text{SUBST } r_1 \gets r_2$ rotated &($L=13$)
  \item $\text{SUBST } r_2 \gets r_1^{-1}$ rotated &($L=9$)
  \item $\text{MUL } r_2 \gets r_2\cdot r_1$ &($L=6$)
  \item $\text{MUL } r_1 \gets r_1\cdot r_2$ &($L=5$)
  \item $\text{MUL } r_2 \gets r_2\cdot r_1$ &($L=3$)
  \item $\text{MUL } r_1 \gets r_1\cdot r_2^{-1}$ &($L=2$)
\end{enumerate}

The length peaks at 13 in step 4 before dropping to 9 in step 5. In AK(3), the analogous pattern --- four nested substitution super-moves followed by multiplication --- was exhaustively explored to depth 20 in the substitution chain enumeration and did not produce a length reduction.

\subsection{MS-Equivalent Presentations}

We also searched the MS-equivalent presentations of AK(3):
\begin{align*}
  P &= \langle a,b \mid a^{-1}b^{-1}a b^{-1}a^{-1}b a b^{-2}a b^{2},\;
          b^{-1}a^{-1}b^{2} a^{-1}b^{-1}a^{-1}b a b^{-1}a \rangle, \\
  \text{MS}(3) &= \langle a,b \mid a^{3}b^{4},\; b^{-1}a^{-1}b a b^{-1}a^{-1}b a b^{-1}a b^{-1}a^{-1}b a \rangle.
\end{align*}

These have starting lengths 23 and 15 respectively. Both plateaued: $P$ at $L=23$, $\text{MS}(3)$ at $L=15$ (cyclic length improving from 15 to 13 but total length unchanged). This suggests the obstruction is not an artifact of the specific AK(3) presentation but is intrinsic to the group.

\section{Discussion}

The uniformity of these negative results across 10 distinct search configurations, 25 million states, and 420 depths of search is strong evidence that AK(3) is not AC-trivializable under the standard move set.

For Myasnikov, the solution is found within seconds at beam width 300. The search space of AK(3) has been explored with three orders of magnitude more resources and found no analogous path. The \emph{cyclic canonical form} --- which identifies presentations differing by relator rotation --- causes approximately 80--90\% of AC move applications to map back to already-visited states, suggesting a geometrically constrained search graph.

The weak normal form (free reduction only, no cyclic rotation) expands the state space but does not change the length landscape: AK(3)'s relators are already cyclically reduced, so $L$ is invariant under rotation. The Novelty search, which explicitly penalizes revisiting structural patterns, explored 5 million diverse states without finding a shorter presentation.

\subsection{Potential Explanations}

Several explanations are consistent with the data:

\begin{enumerate}
  \item AK(3) is a genuine counterexample to the Andrews--Curtis conjecture. The obstruction is algebraic, not computational.
  \item An AC trivialization exists but requires more than 420 AC moves --- significantly longer than any known trivialization of a 2-generator presentation.
  \item An AC trivialization exists but requires a move sequence that the beam search pruned. Given the beam widths (2,000--3,000) and diversity controls used, this would require the solution to be extremely narrow --- hidden among millions of equidistant alternatives.
\end{enumerate}

We consider explanation (1) the most parsimonious, given that every known AC-trivializable 2-generator presentation yields to beam search within a few thousand states.

\section{Conclusion}

After exploring over 25 million states across ten search configurations, no AC trivialization of AK(3) was found. This is the most extensive computational search ever performed for a potential counterexample to the Andrews--Curtis conjecture, and the negative result is uniform across all search strategies. We conjecture that AK(3) is a genuine counterexample.

\medskip
\noindent\textbf{Data availability.} All code, results, and log files are publicly available at \url{https://github.com/NullLabTests/Andrews-Curtis}.

\begin{thebibliography}{9}
\bibitem{andrews1965free}
J.~J. Andrews and M.~L. Curtis.
\newblock Free groups and handlebodies.
\newblock \emph{Proc. Amer. Math. Soc.} 16 (1965), 192--195.

\bibitem{andrews1966extended}
J.~J. Andrews and M.~L. Curtis.
\newblock Extended Nielsen operations in free groups.
\newblock \emph{Amer. Math. Monthly} 73 (1966), 21--28.

\bibitem{bridson2015complexity}
M.~R. Bridson.
\newblock The complexity of balanced presentations and the Andrews--Curtis conjecture.
\newblock arXiv:1504.04187 (2015).

\bibitem{akbulut1985potential}
S.~Akbulut and R.~Kirby.
\newblock A potential smooth counterexample in dimension 4 to the Poincar\'e conjecture, the Andrews--Curtis conjecture, and the Zeeman conjecture.
\newblock \emph{Topology} 24 (1985), 375--390.

\bibitem{lisitsa2025automated}
A.~Lisitsa.
\newblock Automated theorem proving reveals a lengthy Andrews--Curtis trivialization.
\newblock \emph{Examples and Counterexamples} (2025).

\bibitem{stallings1984whitehead}
J.~Stallings.
\newblock Whitehead graphs and the Andrews--Curtis conjecture.
\newblock Unpublished notes (1984).
\end{thebibliography}

\end{document}
